\documentclass[11pt,a4paper]{article}
\usepackage[T1]{fontenc}
\usepackage[utf8]{inputenc}
\usepackage[margin=22mm]{geometry}
\usepackage{lmodern}
\usepackage{amsmath,amssymb}
\usepackage{enumitem}
\usepackage{xcolor}
\usepackage{microtype}
\usepackage{hyperref}
\definecolor{epflred}{HTML}{E3000B}
\hypersetup{colorlinks=true,urlcolor=epflred,linkcolor=epflred,
  pdftitle={ENG-654 Project 6: Offset-Wrist IK through the FANUC CRX Case Study}}
\setlength{\parindent}{0pt}
\setlength{\parskip}{6pt}
\setlist[itemize]{leftmargin=*,itemsep=5pt,topsep=4pt}
\urlstyle{tt}

\begin{document}
{\small\textbf{EPFL \textbar\ ENG-654}\hfill Kinematics-Grounded Motion Planning for Robots}
\vspace{5mm}

{\color{epflred}\Large\bfseries Project 6}

{\LARGE\bfseries Offset-Wrist IK through the FANUC CRX Case Study\par}

\textbf{Format:} Group of two students; simulation-based project.\\
\textbf{Prerequisites:} Rigid transforms, D--H modelling, numerical inverse
kinematics, Jacobians, singularities, and polynomial root finding.

\section*{Motivation}
The FANUC CRX provides a concrete example of an offset-wrist 6R robot
whose inverse kinematics cannot be treated as the PUMA arm--spherical-wrist
decomposition. Studying a richer solution set reveals why IK enumeration and
global connectivity must be examined separately before planning a motion.

\section*{Task}
\textbf{Overall objectives.} Use the supplied CRX-10iA/L IK model to study solution multiplicity,
singularities, and selected connections between IK solutions. Federico Thomas
and Josep M. Porta have developed the robot's closure-polynomial formulation;
the numerical IK implementation will be provided, so the project focuses on
interpreting its solutions.

\begin{itemize}
  \item \textbf{Validate the CRX model.} Inspect the supplied URDF/STL geometry and
document the D--H model, joint signs and offsets, and base/tool frames used by
the IK implementation. Compare forward kinematics from both representations.
Explain geometrically why a spherical-wrist decomposition is unavailable.

  \item \textbf{Understand the polynomial formulation.} Read the CRX case in
\emph{The inverse kinematics of lobster arms}. Describe the distance variable,
closure polynomial, and recovery of joint configurations. Distinguish an
explicit polynomial formulation from a formula for its roots in radicals;
re-deriving the full polynomial is not required.

  \item \textbf{Inspect and validate solution multiplicity.} Reproduce the supplied
sixteen-solution reference case after checking the angle convention. Verify
position and orientation closure for every solution, deduplicate periodic
equivalents, and report joint-limit admissibility separately. Compare two
nearby target poses using the supplied solver.

  \item \textbf{Visualize singularities and IK locations.} Evaluate $\det J$ on two
selected two-joint slices, stating all fixed joint values; choose a
$(q_2,q_3)$ slice for at least one plot. A projected IK whose other joints differ
does not belong to that slice: distinguish projections from on-slice points.
Supplement the plots with full-Jacobian singular values at every IK.

  \item \textbf{Investigate one connectivity question.} Use an instructor-supplied
candidate joint path between two IKs of a common pose. Check endpoint pose
agreement, full-Jacobian regularity, and the closed workspace image of the
path. Explain why the existence of sixteen solutions or equal determinant
signs alone does not establish their connectivity or cuspidality.
\end{itemize}

\newpage
\section*{Specifics and scope}
Use one reference target, two small pose perturbations, two joint-space
slices, and one supplied candidate path. The IK model and candidate path will
be provided. Numerical polynomial solution, full-pose checks, and interpretation
are required; deriving the general 6R IK polynomial or searching the entire
six-dimensional configuration space is outside scope. State the scaling used
for the mixed linear/angular Jacobian.

Check both translation and orientation residuals
for every claimed full-pose IK. Document angle periodicity and limits,
Jacobian conventions, and any scaling used for singular values. State all
fixed coordinates in reduced plots; a two-dimensional slice does not show
the entire 6R singular set. Refine decisive transitions, and distinguish
sampled evidence from a continuous-path guarantee. Collision freedom is
not implied by these kinematic checks.

\section*{Expected outcome}
Understand CRX geometry, the role of closure polynomials, the richer IK
distribution, and the distinction between solution count and connectivity.
Deliver a validated IK table, singularity slices, and an evidence-based analysis
of the candidate path and its workspace image.

Submit reproducible code, model parameters and test data, the required plots,
a concise 5--6-page report, and a short simulation demonstration. Explain
both successful cases and failures; conclusions must follow the numerical
and geometric evidence rather than an expected outcome.

\section*{Provided materials and implementation}
The CRX-10iA/L robot description (.urdf) and link meshes (.stl) will be provided
to the students, together with the numerical IK model, sixteen-solution
reference data, and candidate joint path. Use the same model variant for all
comparisons. The reference joint angles are given in degrees in a D--H
convention. Let $\theta_i$ denote these angles converted to radians, and map
them to the supplied URDF convention using
\[
\begin{aligned}
q_i&=\operatorname{wrap}_{[-\pi,\pi)}(s_i\theta_i+b_i),\\
s&=(1,-1,-1,1,-1,-1),\qquad b=(0,-\pi/2,0,\pi,\pi,0).
\end{aligned}
\]
Check this conversion by comparing all reference forward poses. Wrapped angles
are representatives modulo $2\pi$; check physical limits separately and keep
joint trajectories continuous when validating the candidate path.

Suggested reading: Federico Thomas and Josep M. Porta,
\emph{The inverse kinematics of lobster arms},
Mechanism and Machine Theory 196 (2024), 105630. The paper develops a
distance-based closure-polynomial formulation, including the CRX-10iA/L case.

Use a robotics library or your own tools to verify D--H parameters, inspect
axes and frames, visualize the recovered IKs and singularity slices, and
animate the candidate path. Validate numerical results independently with
forward kinematics and the full Jacobian.

\end{document}
